\documentclass[11pt]{article}
% Shared preamble for the printable mini lessons (1.4 onward).
%   \documentclass[11pt]{article}
%   % Shared preamble for the printable mini lessons (1.4 onward).
%   \documentclass[11pt]{article}
%   % Shared preamble for the printable mini lessons (1.4 onward).
%   \documentclass[11pt]{article}
%   \input{../preamble}
% Build the key with:  pdflatex -jobname=<name>_KEY "\def\iskey{}\input{<name>}"
% Figures shared with the web page live in figures/Unit_X/<name>.tex and are
% pulled in with \fig{Unit_X/<name>}; the standalone package drops their own
% preamble, so their macros come from figures/preamble.tex, loaded here.
\input{../../figures/preamble}
\usepackage[margin=0.75in]{geometry}
\usepackage{array}
\usepackage{enumitem}
\usepackage{amsmath}
\usepackage[helvet]{sfmath}  % math in the same sans font as the text
\usepackage{standalone}
\pagestyle{empty}
\setlength{\parindent}{0pt}
\setlist[enumerate,1]{leftmargin=1.5em,itemsep=1em}
\setlist[enumerate,2]{label=(\alph*),leftmargin=2em,itemsep=2.2em}

\newif\ifkey
\ifdefined\iskey \keytrue \fi
\colorlet{keyred}{red!75!black}
\newcommand{\ans}[1]{\ifkey\leavevmode{\color{keyred}#1}\fi}
% vertical work space on the student copy, less on the key
\newcommand{\work}[1]{\ifkey\vspace{0.6em}\else\vspace{#1}\fi}
\newcommand{\fig}[1]{\input{../../figures/#1}}
% Title line: \lessontitle{Quartiles, IQR, and Box Plots}
\newcommand{\lessontitle}[1]{{\Large\bfseries #1\ifkey\ \textcolor{keyred}{KEY}\fi}\hfill Name \underline{\hspace{2in}}\par\vspace{0.6em}}
\newcommand{\partner}{\par\vspace{0.4em}Work with a partner. Compare answers when you finish each problem.\par}

% Build the key with:  pdflatex -jobname=<name>_KEY "\def\iskey{}\input{<name>}"
% Figures shared with the web page live in figures/Unit_X/<name>.tex and are
% pulled in with \fig{Unit_X/<name>}; the standalone package drops their own
% preamble, so their macros come from figures/preamble.tex, loaded here.
% Shared preamble for web figures.  Each figure is a standalone .tex that
% \input's this file; scripts/make-figures.sh turns it into an SVG.
\usepackage[T1]{fontenc}
\usepackage{helvet}
\renewcommand{\familydefault}{\sfdefault}
\usepackage{tikz}
\usetikzlibrary{decorations.pathreplacing}
% axis numbers in the text font, not math italic
\pgfkeys{/pgf/number format/.cd,assume math mode=true,1000 sep={}}

% dot plot: \dotplot{min}{max}{value/count,...}{axis label}
\newcommand{\dotplot}[4]{%
\begin{tikzpicture}[x=0.9cm,y=0.33cm]
  \draw[thick] (#1-0.5,0) -- (#2+0.5,0);
  \foreach \t in {#1,...,#2} {\draw (\t,0.25) -- (\t,-0.25) node[below] {\t};}
  \foreach \v/\c in {#3} {\ifnum\c>0 \foreach \k in {1,...,\c} {\fill (\v,\k*1.15) circle (0.16cm);}\fi}
  \node at ({(#1+#2)/2},-2.3) {#4};
\end{tikzpicture}}

% shape sketch histogram (12 bins): \shapehist{heights}{label}
\newcommand{\shapehist}[2]{%
  \foreach \h [count=\i from 0] in {#1} \draw[fill=gray!30] (\i,0) rectangle (\i+1,\h);
  \draw (-0.3,0) -- (12.3,0);
  \node[font=\bfseries] at (6,-2.6) {#2};}
\newcommand{\meanmark}[1]{\fill (#1,0) -- ++(-0.3,-0.75) -- ++(0.6,0) -- cycle;}
\newcommand{\medmark}[2]{\draw[very thick,dashed] (#1,0) -- (#1,#2);}

% number line with minor ticks: \nline{min}{max}{label step}{axis label}
\newcommand{\nline}[4]{%
  \draw[thick] (#1-0.5,0) -- (#2+0.5,0);
  \foreach \t in {#1,...,#2} \draw (\t,0.12) -- (\t,-0.12);
  \pgfmathtruncatemacro{\nxt}{#1+#3}
  \foreach \t in {#1,\nxt,...,#2} \draw (\t,0.2) -- (\t,-0.2) node[below] {\t};
  \node at ({(#1+#2)/2},-0.95) {#4};}
% box plot above the line: \bplot{min}{Q1}{med}{Q3}{max}
\newcommand{\bplot}[5]{%
  \draw[thick] (#1,1.0) -- (#2,1.0);
  \draw[thick] (#4,1.0) -- (#5,1.0);
  \draw[thick] (#1,0.75) -- (#1,1.25);
  \draw[thick] (#5,0.75) -- (#5,1.25);
  \draw[thick,fill=gray!20] (#2,0.55) rectangle (#4,1.45);
  \draw[very thick] (#3,0.55) -- (#3,1.45);}

% ---- added for 1.4 onward ----
% axis on y=0: \axis{min}{max}{label step}{tick step}{axis label}
\newcommand{\axis}[5]{%
  \draw[thick] (#1,0) -- (#2,0);
  \pgfmathsetmacro{\axT}{#1+#4}\foreach \t in {#1,\axT,...,#2} \draw (\t,0.1) -- (\t,-0.1);
  \pgfmathsetmacro{\axL}{#1+#3}\foreach \t in {#1,\axL,...,#2}
    \draw (\t,0.18) -- (\t,-0.18) node[below] {\pgfmathprintnumber{\t}};
  \node at ({(#1+#2)/2},-0.95) {#5};}
% box plot centred at height y: \bplotat{y}{min}{Q1}{med}{Q3}{max}
\newcommand{\bplotat}[6]{%
  \draw[thick] (#2,#1) -- (#3,#1);
  \draw[thick] (#5,#1) -- (#6,#1);
  \draw[thick] (#2,#1-0.22) -- (#2,#1+0.22);
  \draw[thick] (#6,#1-0.22) -- (#6,#1+0.22);
  \draw[thick,fill=gray!20] (#3,#1-0.42) rectangle (#5,#1+0.42);
  \draw[very thick] (#4,#1-0.42) -- (#4,#1+0.42);}
% an outlier dot: \outlier{x}{y}
\newcommand{\outlier}[2]{\fill (#1,#2) circle (2.4pt);}

\usepackage[margin=0.75in]{geometry}
\usepackage{array}
\usepackage{enumitem}
\usepackage{amsmath}
\usepackage[helvet]{sfmath}  % math in the same sans font as the text
\usepackage{standalone}
\pagestyle{empty}
\setlength{\parindent}{0pt}
\setlist[enumerate,1]{leftmargin=1.5em,itemsep=1em}
\setlist[enumerate,2]{label=(\alph*),leftmargin=2em,itemsep=2.2em}

\newif\ifkey
\ifdefined\iskey \keytrue \fi
\colorlet{keyred}{red!75!black}
\newcommand{\ans}[1]{\ifkey\leavevmode{\color{keyred}#1}\fi}
% vertical work space on the student copy, less on the key
\newcommand{\work}[1]{\ifkey\vspace{0.6em}\else\vspace{#1}\fi}
\newcommand{\fig}[1]{\input{../../figures/#1}}
% Title line: \lessontitle{Quartiles, IQR, and Box Plots}
\newcommand{\lessontitle}[1]{{\Large\bfseries #1\ifkey\ \textcolor{keyred}{KEY}\fi}\hfill Name \underline{\hspace{2in}}\par\vspace{0.6em}}
\newcommand{\partner}{\par\vspace{0.4em}Work with a partner. Compare answers when you finish each problem.\par}

% Build the key with:  pdflatex -jobname=<name>_KEY "\def\iskey{}\input{<name>}"
% Figures shared with the web page live in figures/Unit_X/<name>.tex and are
% pulled in with \fig{Unit_X/<name>}; the standalone package drops their own
% preamble, so their macros come from figures/preamble.tex, loaded here.
% Shared preamble for web figures.  Each figure is a standalone .tex that
% \input's this file; scripts/make-figures.sh turns it into an SVG.
\usepackage[T1]{fontenc}
\usepackage{helvet}
\renewcommand{\familydefault}{\sfdefault}
\usepackage{tikz}
\usetikzlibrary{decorations.pathreplacing}
% axis numbers in the text font, not math italic
\pgfkeys{/pgf/number format/.cd,assume math mode=true,1000 sep={}}

% dot plot: \dotplot{min}{max}{value/count,...}{axis label}
\newcommand{\dotplot}[4]{%
\begin{tikzpicture}[x=0.9cm,y=0.33cm]
  \draw[thick] (#1-0.5,0) -- (#2+0.5,0);
  \foreach \t in {#1,...,#2} {\draw (\t,0.25) -- (\t,-0.25) node[below] {\t};}
  \foreach \v/\c in {#3} {\ifnum\c>0 \foreach \k in {1,...,\c} {\fill (\v,\k*1.15) circle (0.16cm);}\fi}
  \node at ({(#1+#2)/2},-2.3) {#4};
\end{tikzpicture}}

% shape sketch histogram (12 bins): \shapehist{heights}{label}
\newcommand{\shapehist}[2]{%
  \foreach \h [count=\i from 0] in {#1} \draw[fill=gray!30] (\i,0) rectangle (\i+1,\h);
  \draw (-0.3,0) -- (12.3,0);
  \node[font=\bfseries] at (6,-2.6) {#2};}
\newcommand{\meanmark}[1]{\fill (#1,0) -- ++(-0.3,-0.75) -- ++(0.6,0) -- cycle;}
\newcommand{\medmark}[2]{\draw[very thick,dashed] (#1,0) -- (#1,#2);}

% number line with minor ticks: \nline{min}{max}{label step}{axis label}
\newcommand{\nline}[4]{%
  \draw[thick] (#1-0.5,0) -- (#2+0.5,0);
  \foreach \t in {#1,...,#2} \draw (\t,0.12) -- (\t,-0.12);
  \pgfmathtruncatemacro{\nxt}{#1+#3}
  \foreach \t in {#1,\nxt,...,#2} \draw (\t,0.2) -- (\t,-0.2) node[below] {\t};
  \node at ({(#1+#2)/2},-0.95) {#4};}
% box plot above the line: \bplot{min}{Q1}{med}{Q3}{max}
\newcommand{\bplot}[5]{%
  \draw[thick] (#1,1.0) -- (#2,1.0);
  \draw[thick] (#4,1.0) -- (#5,1.0);
  \draw[thick] (#1,0.75) -- (#1,1.25);
  \draw[thick] (#5,0.75) -- (#5,1.25);
  \draw[thick,fill=gray!20] (#2,0.55) rectangle (#4,1.45);
  \draw[very thick] (#3,0.55) -- (#3,1.45);}

% ---- added for 1.4 onward ----
% axis on y=0: \axis{min}{max}{label step}{tick step}{axis label}
\newcommand{\axis}[5]{%
  \draw[thick] (#1,0) -- (#2,0);
  \pgfmathsetmacro{\axT}{#1+#4}\foreach \t in {#1,\axT,...,#2} \draw (\t,0.1) -- (\t,-0.1);
  \pgfmathsetmacro{\axL}{#1+#3}\foreach \t in {#1,\axL,...,#2}
    \draw (\t,0.18) -- (\t,-0.18) node[below] {\pgfmathprintnumber{\t}};
  \node at ({(#1+#2)/2},-0.95) {#5};}
% box plot centred at height y: \bplotat{y}{min}{Q1}{med}{Q3}{max}
\newcommand{\bplotat}[6]{%
  \draw[thick] (#2,#1) -- (#3,#1);
  \draw[thick] (#5,#1) -- (#6,#1);
  \draw[thick] (#2,#1-0.22) -- (#2,#1+0.22);
  \draw[thick] (#6,#1-0.22) -- (#6,#1+0.22);
  \draw[thick,fill=gray!20] (#3,#1-0.42) rectangle (#5,#1+0.42);
  \draw[very thick] (#4,#1-0.42) -- (#4,#1+0.42);}
% an outlier dot: \outlier{x}{y}
\newcommand{\outlier}[2]{\fill (#1,#2) circle (2.4pt);}

\usepackage[margin=0.75in]{geometry}
\usepackage{array}
\usepackage{enumitem}
\usepackage{amsmath}
\usepackage[helvet]{sfmath}  % math in the same sans font as the text
\usepackage{standalone}
\pagestyle{empty}
\setlength{\parindent}{0pt}
\setlist[enumerate,1]{leftmargin=1.5em,itemsep=1em}
\setlist[enumerate,2]{label=(\alph*),leftmargin=2em,itemsep=2.2em}

\newif\ifkey
\ifdefined\iskey \keytrue \fi
\colorlet{keyred}{red!75!black}
\newcommand{\ans}[1]{\ifkey\leavevmode{\color{keyred}#1}\fi}
% vertical work space on the student copy, less on the key
\newcommand{\work}[1]{\ifkey\vspace{0.6em}\else\vspace{#1}\fi}
\newcommand{\fig}[1]{\input{../../figures/#1}}
% Title line: \lessontitle{Quartiles, IQR, and Box Plots}
\newcommand{\lessontitle}[1]{{\Large\bfseries #1\ifkey\ \textcolor{keyred}{KEY}\fi}\hfill Name \underline{\hspace{2in}}\par\vspace{0.6em}}
\newcommand{\partner}{\par\vspace{0.4em}Work with a partner. Compare answers when you finish each problem.\par}


\begin{document}

\lessontitle{Experiments}

Does a new pill stop headaches? We cannot just compare people who chose to take it with people who did not.
They may be different in other ways. An \textbf{experiment} gives the treatment on purpose and compares fairly.

\vspace{0.5em}
\renewcommand{\arraystretch}{1.3}
\begin{tabular}{|>{\raggedright\arraybackslash}p{1.7in}|>{\raggedright\arraybackslash}p{4.8in}|}\hline
\textbf{Word} & \textbf{In the headache experiment}\\ \hline
Subjects (experimental units) & The 40 volunteers with headaches\\ \hline
Explanatory variable & What we change: the type of pill\\ \hline
Treatments & The things we compare: new pill, placebo\\ \hline
Response variable & What we measure: headache relief\\ \hline
Control group & The group that does not get the new treatment (the placebo group)\\ \hline
Placebo & A fake treatment that looks real, like a sugar pill\\ \hline
Double-blind & Neither the subjects nor the people measuring know who got which pill\\ \hline
\end{tabular}

\begin{center}\fig{Unit_5/5.3_design}\end{center}

\textbf{Principles of a good experiment}

\vspace{0.2em}
\begin{tabular}{@{}lp{5.25in}}
\textbf{Control} & Compare to a control group. Keep everything else the same for both groups.\\
\textbf{Random assignment} & Use chance to put subjects into groups. Then the groups are alike except for the treatment.\\
\textbf{Replication} & Use enough subjects that chance differences between people even out.\\
\end{tabular}

\vspace{0.6em}
\textbf{Why a placebo?} People often feel better just because they \textit{think} they got medicine (the \textbf{placebo effect}).
With a placebo group, both groups think they are being treated.

\vspace{0.6em}
\textbf{Matched pairs design.} Each subject gets \textbf{both} treatments, in a random order. Or pair up very similar subjects, like twins.

\vspace{0.2em}
\textbf{Example:} Which running shoe is faster? Each runner runs once in each shoe. A coin flip decides which shoe goes first.
Comparing each runner to themselves removes the differences between runners.

\partner

% ---------------- PROBLEMS ----------------
\newpage
\begin{enumerate}

\item 60 tomato plants are randomly split into two groups of 30. One group gets a new fertilizer. The other gets only water. After 6 weeks, the weight of the tomatoes is measured.
\begin{enumerate}[itemsep=1.4em]
  \item What are the subjects? \ans{\quad The 60 tomato plants}
  \item What is the explanatory variable? \ans{\quad Fertilizer or water (the treatment)}
  \item What is the response variable? \ans{\quad The weight of the tomatoes}
  \item What is the control group? \ans{\quad The 30 plants that get only water}
  \item Why split the plants \textbf{randomly}? \ans{\quad So the groups are alike. Then a difference comes from the fertilizer, not from which plants were stronger.}
\end{enumerate}
\work{2em}

\item In the headache experiment, the study is double-blind.
\begin{enumerate}[itemsep=1em]
  \item Why should the subjects not know which pill they got?

  \ans{If they know they got the real pill, they might feel better just from believing it (the placebo effect).}
  \work{2.5em}
  \item Why should the nurses who measure headache relief not know?

  \ans{They might, without meaning to, rate the new-pill group better.}
\end{enumerate}

\newpage
\item Does listening to music help students memorize a list of words? You have 50 student volunteers.
Describe a completely randomized experiment. Name the treatments, how you assign them, and what you measure.

\ans{Treatments: music, no music. Number the students 1 to 50 and use a random number generator to put 25 in each group. Both groups study the same list for the same time in the same room. Measure how many words each student remembers. Compare the two groups.}
\work{7em}

\item A company tests two hand lotions. Each of 20 volunteers puts Lotion A on one hand and Lotion B on the other. A coin flip decides which hand gets A. After a week, a doctor rates the dryness of each hand.
\begin{enumerate}[itemsep=1em]
  \item What kind of design is this? \ans{\quad Matched pairs}
  \item Why flip a coin to pick the hand? \ans{\quad Most people use one hand more. The coin keeps that from favoring one lotion.}
  \work{1.5em}
  \item Why is this design better than giving 10 people Lotion A and 10 people Lotion B?

  \ans{Each person is compared to themselves, so differences in skin from person to person do not matter.}
\end{enumerate}
\work{2em}

\item A teacher tries a new reading program with her 1st-hour class and the old program with her 6th-hour class. The 1st-hour class scores higher. Can she conclude the new program works better? Explain.

\ans{No. She did not use random assignment. The classes could differ in other ways (time of day, which students take 1st hour). Those differences could cause the higher scores.}

\end{enumerate}

\end{document}
